\begin{enumerate}
	\item STVU est un rectangle, donc ses côtés opposés sont de même longueur : ST = UV$=27$~m.

	Les points S, T et M sont alignés dans cet ordre, donc :
	\[\text{TM}=\text{SM} - \text{ST} =36-27=9~\text{m}.\]

	\item STVU étant un rectangle, TV = SU $=12$~m. Le triangle TMV est rectangle en T. D'après le théorème de Pythagore :
	\[\text{MV}^{2}=\text{TM}^{2}+\text{TV}^{2} \]
	\[\text{MV}^{2}=9^{2}+12^{2}\]
	\[\text{MV}^{2}=81+144\]
	\[\text{MV}^{2}=225.\]
	Donc MV$ = \sqrt{225}=15$~m.

	\item Le périmètre du terrain SMVU est
\[\text{SM}+\text{MV}+\text{VU}+\text{US} = 36 + 15 + 27 + 12 = 90~\text{m}.\]
	
En retirant l'ouverture de 3~m pour le portail : $90-3=87$.

	Il faut prévoir 87~m de grillage.

	\item STVU est un rectangle de longueur UV$=27$~m et de largeur SU$=12$~m :

\[\mathcal{A}= \text{UV}\times \text{SU}=27\times 12 = 324~\text{m}^2.\]

	\item Calculons l'aire d'un rouleau : $4\times 2= 8~$(m$^2$).

	Nombre de rouleaux nécessaires : $324\div 8=40,5$ ; il faut donc acheter 41 rouleaux.

	$41\times 86=\np{3526}$~(\euro).
	
Le coût pour l'achat de 41 rouleaux de gazon est de $\np{3526}$~\euro.

	Comme $\np{3526}$ est supérieur à $\np{3500}$, le budget de \np{3500}~\euro{} n'est pas suffisant.
\end{enumerate}
